% =============================================================
%  ch22_merkmalserkennung.tex  --  Merkmalserkennung ohne KI-Training
% =============================================================

\chapter{Merkmalserkennung ohne KI-Training}
\label{ch:merkmalserkennung}

\section{Detektoren im Vergleich: SIFT, SURF, KAZE, AKAZE, ORB, BRISK}
Bevor wir Code schreiben, lohnt der nüchterne Vergleich der klassischen
Detektor-Deskriptor-Paare. Sie durchlaufen alle dieselben Stufen --
Merkmalsdetektion und -beschreibung, Matching, Ausreißerverwerfung
(Lowe-Ratio + RANSAC) und Schätzung der Transformation -- unterscheiden sich
aber stark in Genauigkeit und Geschwindigkeit \cite{tareen2018comparative}:
\begin{itemize}
\item \textbf{SIFT} und \textbf{KAZE/AKAZE} liefern die \emph{genauesten},
  invariantesten Matches (Skala, Rotation, Blickwinkel), sind aber rechenintensiver.
\item \textbf{ORB} und \textbf{BRISK} verwenden \emph{binäre} Deskriptoren mit
  Hamming-Matching und sind dadurch am \emph{schnellsten} und sparsamsten --
  ideal für Echtzeit auf schwacher Hardware, zum Preis etwas geringerer
  Robustheit. ORB findet zudem typischerweise die meisten Merkmale pro Bild.
\end{itemize}

\begin{figure}[htbp]
  \centering
  \includegraphics[width=0.85\textwidth]{image_detection}
  \caption{Merkmalsdetektion, Matching und Bildregistrierung derselben Szene mit
  SIFT, SURF, KAZE, AKAZE, ORB und BRISK (Zeilen). Die Detektoren unterscheiden
  sich deutlich in Anzahl und Verteilung der gefundenen Merkmale sowie in der
  Matching-Qualität; SIFT/AKAZE gelten als am genauesten, ORB/BRISK als am
  effizientesten. Nach \cite{tareen2018comparative}.}
  \label{fig:image_detection}
\end{figure}

\noindent \textbf{Faustregel für dieses Projekt:} ORB als Standardwahl für
Echtzeit-Matching und visuelle Odometrie, SIFT/AKAZE, wenn es bei der
CAD-Posenschätzung auf maximale Genauigkeit ankommt. Der folgende Code zeigt
die ORB-Variante exemplarisch.

\begin{lstlisting}[style=py,caption={ORB-Erkennung und Matching (klassisch, kein Training nötig)}]
import cv2

orb = cv2.ORB_create(nfeatures=1500)
# Referenzansicht des Bauteils (z.B. aus CAD-Render) einmalig beschreiben:
kp_ref, des_ref = orb.detectAndCompute(ref_gray, None)
matcher = cv2.BFMatcher(cv2.NORM_HAMMING)

def detect(frame_gray):
    kp, des = orb.detectAndCompute(frame_gray, None)
    if des is None:
        return []
    raw  = matcher.knnMatch(des_ref, des, k=2)
    good = [m for m, n in raw if m.distance < 0.75 * n.distance]  # Lowe
    return kp, good

# 6D-Pose ueber 2D-3D-Korrespondenzen (CAD-Punkte <-> Bildpunkte):
# ok, rvec, tvec = cv2.solvePnPRansac(obj_pts_3d, img_pts_2d,
#                                      K, distCoeffs)
\end{lstlisting}

\noindent Für \textbf{Bewegungserkennung} (Projekt~3) genügt
Hintergrundsubtraktion:
\begin{lstlisting}[style=py,caption={Bewegungserkennung ohne Merkmale}]
bg = cv2.createBackgroundSubtractorMOG2(detectShadows=False)
mask = bg.apply(frame)                      # bewegte Pixel
contours, _ = cv2.findContours(
    mask, cv2.RETR_EXTERNAL, cv2.CHAIN_APPROX_SIMPLE)
movers = [c for c in contours if cv2.contourArea(c) > 500]
\end{lstlisting}
